% Les trois piles de protocoles du web : ce que chaque version de HTTP
% délègue à la couche inférieure.
\documentclass[border=6pt]{standalone}
\usepackage{sd}
\begin{document}
\begin{tikzpicture}[papier, x=1mm, y=1mm,
    c/.style={bloc=#1, minimum width=40mm, text width=37mm, minimum height=9mm}]
  \node[titrec, anchor=south] at (0,6) {1997 (RFC 2068)};
  \node[titrec, anchor=south] at (50,6) {2015 (RFC 7540)};
  \node[titrec, anchor=south] at (100,6) {2022 (RFC 9114)};
  % HTTP/1.1
  \node[c=Pbleu] at (0,0) {HTTP/1.1\\[-1pt]{\scriptsize texte, une requête à la fois}};
  \node[c=Plavande] at (0,-11) {TLS};
  \node[c=Porange, minimum height=20mm] at (0,-26.5) {TCP\\[-1pt]{\scriptsize fiabilité, ordre, congestion}};
  \node[c=Pfond] at (0,-42) {IP};
  % HTTP/2
  \node[c=Pbleu] at (50,0) {HTTP/2\\[-1pt]{\scriptsize trames binaires, flux multiplexés, HPACK}};
  \node[c=Plavande] at (50,-11) {TLS};
  \node[c=Porange, minimum height=20mm] at (50,-26.5) {TCP\\[-1pt]{\scriptsize un seul flux d'octets ordonné}};
  \node[c=Pfond] at (50,-42) {IP};
  % HTTP/3
  \node[c=Pbleu] at (100,0) {HTTP/3\\[-1pt]{\scriptsize trames, QPACK}};
  \node[c=Pvert, minimum height=20mm] at (100,-15.5) {QUIC\\[-1pt]{\scriptsize flux indépendants, fiabilité, congestion, TLS 1.3 intégré}};
  \node[c=Porange] at (100,-31) {UDP};
  \node[c=Pfond] at (100,-42) {IP};
  \node[note, anchor=north] at (0,-48) {noyau : TCP ; bibliothèque : TLS};
  \node[note, anchor=north] at (50,-48) {noyau : TCP ; bibliothèque : TLS};
  \node[note, anchor=north] at (100,-48) {noyau : UDP ; application : QUIC};
\end{tikzpicture}
\end{document}
